\section{A full-power reduction for binary Heisenberg factors}
\label{sec:binary-heisenberg-power-profile}

\paragraph{Status and provenance.}
D-TD59, question Q2fo, is separate from the frozen concavity
theorem D-TD57.  The author developed the cyclic low-rank-power argument
while investigating the remaining short projected tuples.  The parent
independently reconstructed its rank, packing and numerical inequalities
before the full proof below was written.
The complete written proof and checker have passed the author's whole
reading, and the exact supplementary certificate passed its first execution.
Complete root and independent whole-proof reviews also pass.
No complete bound on the binary Heisenberg exponent is claimed.

\paragraph{Actual source, without a whole-image floor.}
Let \(C\to\mathbf P^1\) be a connected finite \(J\)-Galois cover with
faithful irreducible complex coefficient representation on \(V\).
Let its actual generating product-one inertia tuple \(T_1,\ldots,T_r\)
satisfy
\[
                e=\dim V,\qquad
                \sum_{i=1}^r\operatorname{rank}(T_i-I)=2e+6.
 \tag{BHP.1}
\]
Suppose \(V=W\otimes U\) is preserved factorwise, with actual finite
projective image \(P\subset\operatorname{PGL}(W)\).  Assume \(P\) is
primitive and has a normal projective Heisenberg subgroup
\[
 Q\simeq\mathbf F_2^{\,2m},\qquad
 [L_v,L_w]=(-1)^{\langle v,w\rangle}I,\qquad
             a=\dim W=2^m,
 \tag{BHP.2}
\]
acting irreducibly on \(W\), with nondegenerate alternating pairing.
Let \(y_i\in P\) be the actual projected inertia, \(g_i\in\operatorname{Sp}(Q)\)
its conjugation action on the label space, and
\[
                  k=\#\{i:y_i\ne1\}.
 \tag{BHP.3}
\]
There is no hypothesis \(r\ge12\), no lower bound for the support of
the entire original image, and no lower bound for partner-only marks.
All scalar choices and the arbitrary whole \(U\) are retained.

\begin{theorem}[D-TD59]
In the above actual source, if \(m>80\), then
\[
                              3\le k\le4.
 \tag{BHP.4}
\]
For the local estimate used in the proof, if a finite-projective-order
normalizer \(A\) induces \(g\in\operatorname{Sp}(Q)\), \(m>10\), and
\(\operatorname{rank}(g-I_Q)>m+10\), then
\[
                         \frac{\nu_W(A)}a\ge\frac{31}{48}.
 \tag{BHP.5}
\]
This is an all-target bound and survives tensoring with an arbitrary
finite-order partner.  The theorem does not exclude the remaining
three- and four-mark projected tuples.
\end{theorem}

\begin{proof}
\emph{Step 1: two rank identities and an exact cyclic packing.}
For invertible linear maps \(x,y\) on a vector space,
\[
 \operatorname{rank}(xy-I)
       \le\operatorname{rank}(x-I)+\operatorname{rank}(y-I),
 \qquad
 \operatorname{rank}(x^{-1}-I)=\operatorname{rank}(x-I).
 \tag{BHP.6}
\]
The first follows from \(xy-I=x(y-I)+(x-I)\), and the second
from \(x^{-1}-I=-x^{-1}(x-I)\).
In characteristic two, \(g^2-I=(g-I)^2\).  Applying the elementary
rank-of-a-product inequality on the \(2m\)-dimensional label space
therefore gives
\[
 \operatorname{rank}(g^2-I)
                  \ge2\operatorname{rank}(g-I)-2m.
 \tag{BHP.7}
\]
One direct justification is
\(\dim\ker N^2\le2\dim\ker N\) for \(N=g-I\): the map \(N\)
from \(\ker N^2\) into \(\ker N\) has kernel \(\ker N\).

Suppose \(\operatorname{rank}(g-I)>m+10\) and \(m>10\).
Then
\[
       \operatorname{rank}(g-I)>20,\qquad
       \operatorname{rank}(g^2-I)>20.
 \tag{BHP.8}
\]
For any positive integer \(L\) with \(g^L=I\), define the actual
cyclic set of low-rank powers
\[
       \mathcal S=\{j\in\mathbf Z/L\mathbf Z:
                       \operatorname{rank}(g^j-I)\le10\}.
 \tag{BHP.9}
\]
If \(s,t\in\mathcal S\), (BHP.6) gives
\(\operatorname{rank}(g^{s-t}-I)\le20\).
Thus \(s-t\) cannot be congruent to \(1,-1,2,-2\) modulo \(L\),
by (BHP.8).  In particular \(L\ge3\), and the three translates
\(\mathcal S,\mathcal S+1,\mathcal S+2\) are pairwise disjoint.
This proves the exact all-period inequality
\[
                            \frac{|\mathcal S|}{L}\le\frac13.
 \tag{BHP.10}
\]
No claim that \(\mathcal S\) is a subgroup is used.

\emph{Step 2: the full-period trace and every scalar target.}
For a unitary normalizer \(A\), TD05 proves
\[
 \frac{|\operatorname{tr}A^j|}{a}
 \le 2^{-\operatorname{rank}(g^j-I)/2}.
 \tag{BHP.11}
\]
This holds also when \(g^j=I\): the trace can then be zero or have
absolute value \(a\).  In particular projected identity powers are
retained, rather than dropped.

The projective group generated by \(Q\) and \([A]\) is finite:
its normal subgroup \(Q\) is finite and its quotient is generated by
the finite-order coset of \([A]\).  Its special-linear inverse image
is finite, so it admits a common invariant positive Hermitian form.
First scale \(A\) so that \(A^L=I\) for some finite \(L\); in this
form the scaled \(A\) is unitary, as in the TD05 interface.
For \(j\notin\mathcal S\), the right side of (BHP.11) is at most
\(1/32\).  For \(j\in\mathcal S\) use the bound one.
Fourier averaging at any eigenvalue \(\lambda\) gives
\[
 \begin{split}
 \frac{\dim\ker(A-\lambda I)}a
 &\le\frac1L\sum_{j=0}^{L-1}\frac{|\operatorname{tr}A^j|}{a}\\
 &\le\frac{|\mathcal S|}{L}
                +\left(1-\frac{|\mathcal S|}{L}\right)\frac1{32}
 \le\frac13+\frac{2}{3\cdot32}
 =\frac{17}{48}.
 \end{split}
 \tag{BHP.12}
\]
This proves (BHP.5).  Scaling back leaves the largest eigenspace
unchanged.  The calculation does not replace any original inertia
period by its label period.

For clarity in the actual source, let \(L_i\) be the full order of
\(T_i\).  Writing \(T_i=A_i\otimes B_i\), its identity \(L_i\)-th power
makes the factor \(L_i\)-th powers reciprocal scalars.  Reciprocal
rescaling makes \(A_i^{L_i}=B_i^{L_i}=I\), after which they may be
unitarized.  Apply (BHP.9)--(BHP.12) with this original \(L_i\).
Every original scalar or projected-identity power is included.
Alternatively, diagonalizing \(B_i\) shows directly that each of its
eigenspaces contributes at most the largest \(A_i\)-eigenspace times
its dimension to any chosen eigenspace of \(T_i\).
Thus all the local projective support bounds used below apply to
the actual \(\operatorname{rank}(T_i-I)/e\), with the entire partner.

\emph{Step 3: the three local rank ranges.}
For any nonidentity projective image \(y_i\), put
\(s_i=\operatorname{rank}(g_i-I_Q)\).
The sharp quarter floor of TD08 and its trace-sensitive bound,
together with (BHP.5), give
\[
 \begin{array}{c|c}
 \text{label rank range}&
 \operatorname{rank}(T_i-I)/e\text{ is at least}\\ \hline
 s_i\le10&1/4\\
 10<s_i\le m+10&31/64\\
 s_i>m+10&31/48.
 \end{array}
 \tag{BHP.13}
\]
The middle row uses
\(\nu(A_i)/a\ge(1-2^{-s_i/2})/2\ge31/64\).
The first row includes nonidentity Weyl elements with \(s_i=0\);
no nonidentity projective image is mistaken for a scalar.

\emph{Step 4: the actual rank sum and the short-tuple contradiction.}
The projected tuple generates the noncyclic group \(P\), which
contains \(Q\).  At most two nonidentity product-one generators
would generate a cyclic group, so \(k\ge3\).
The quarter floor and (BHP.1) give
\[
                 k/4\le2+6/e,\qquad
                 e\ge2^m\ge2^{81}>24,
                 \qquad k\le8.
 \tag{BHP.14}
\]
There is no lower bound of twelve on this projected count.

Let \(H=\langle g_i\rangle\).  Primitivity of \(P\) implies \(Q^H=0\):
a nonzero fixed label would give a Weyl operator with two equal
eigenspaces permuted by \(P\), contradicting primitivity.
The exact actual Scott interface proved in TD10 therefore gives
\(\sum_i s_i\ge4m\).  Its application uses the product-one generating
tuple and symplectic self-duality, not an assumption of absolute
irreducibility of the label module.

Let \(\ell\), \(h\), and \(s\) count the large, middle, and small
rank ranges among the \(k\) nonidentity projected entries.
Since every label rank is at most \(2m\),
\[
 \begin{split}
 4m\le\sum_i s_i
 &\le2m\ell+(m+10)h+10s\\
 &\le m(2\ell+h)+10k.
 \end{split}
 \tag{BHP.15}
\]
If \(2\ell+h\le3\), the last expression is at most \(3m+80<4m\).
Hence \(2\ell+h\ge4\).

Suppose now \(k\ge5\).  From (BHP.13) the projected contribution
to the original normalized budget is at least
\[
 \frac{s}{4}+\frac{31h}{64}+\frac{31\ell}{48}
       =\frac{k}{4}+\frac{15h}{64}+\frac{19\ell}{48}
       \ge\frac{49}{24}.
 \tag{BHP.16}
\]
The last minimum can be checked without relaxing any integral count.
If \(\ell\ge2\), the displayed expression is at least
\(5/4+38/48=49/24\).  If \(\ell=1\), then \(h\ge2\) and it is
at least \(203/96>49/24\).  If \(\ell=0\), then \(h\ge4\) and it
is at least \(35/16>49/24\).

Discarding only nonnegative costs at projected identity entries,
(BHP.1) and (BHP.16) imply \(2+6/e\ge49/24\), or \(e\le144\).
This contradicts \(e\ge2^{81}\).  Therefore \(k\le4\), as required.
\end{proof}

\paragraph{Exact remaining boundary.}
This eliminates the binary \(k=5\) branch at \(m>80\), and also
excludes all larger projected branch counts there, without a
whole-original support floor or an original twelve-mark premise.
It leaves \(k=3,4\) and every \(m\le80\) outside its conclusion.
All original scalar phases, arbitrary partners and full inertia
periods have been included.
No realization of a projected tuple as a sixfold source is asserted.
In particular the result does not bound the exponent in either
remaining short branch, classify their label groups, decide a
field or polarization, or supply a detecting algebraic cycle.
